\rok{Фебруар 2026.}{B}

Други поправни тест из Алгоритама и структура података

Први практични тест одржан 19.02.2026.

\zadatak{1}{Bubble Sort}
Написати методу \texttt{Sort} алгоритма \texttt{BubbleSort} за сортирање целобројног низа дужине $n$.

\textbf{Додатне напомене за решавање задатка:}
\begin{itemize}
	\item Улаз је несортиран низ целих бројева.
	\item Излаз је сортирани низ целих бројева.
	\item У template-у се налази класа \texttt{RandomArrayGenerator} коју треба искористити за генерисање низа случајних бројева.
	\item У оквиру класе \texttt{BubbleSort} написати имплементацију за методу:
	\begin{Verbatim}[fontsize=\small]
public static void sort(long[] arr)
	\end{Verbatim}
	\item Обезбедити код у Main методи за тестирање алгоритма.
\end{itemize}

\zadatak{2}{Сума елемената након првог појављивања вредности веће од просека}
Имплементирати методу која у првом пролазу израчунава просечну вредност свих елемената у листи. У другом пролазу треба пронаћи први чвор чија је вредност строго већа од тог просека и вратити суму свих елемената који се налазе после тог чвора (не укључујући њега).

\textbf{Додатни захтеви за решавање задатка:}
\begin{itemize}
	\item Задатак решити у линеарном времену $O(n)$.
	\item Ако ниједан елемент није већи од просека или је тај елемент последњи у листи, вратити 0.
	\item Захтеван потпис методе:
	\begin{Verbatim}[fontsize=\small]
public int sumAfterFirstAboveAverage()
	\end{Verbatim}
\end{itemize}

\textbf{Пример:}
\begin{itemize}
	\item \textbf{Улаз:} $1 \rightarrow 2 \rightarrow 3 \rightarrow 10$ (просек = 4)
	\item \textbf{Излаз:} $0$ (први већи од просека је 10. Пошто после њега нема ничега, output је 0)
	\item \textbf{Улаз:} $10 \rightarrow 2 \rightarrow 1 \rightarrow 3$ (просек = 4)
	\item \textbf{Излаз:} $6$ (први већи од просека је 10. Output је $2 + 1 + 3 = 6$)
\end{itemize}

\zadatak{3}{Селективни одјеци (The Selective Echo)}
Процесираш низ звучних сигнала који се чувају у дигиталном баферу. Сваки сигнал је представљен целим бројем. Због интерференције, нови сигнали могу поништити или очистити претходне на основу своје парности:

\begin{itemize}
	\item \textbf{Поништавање парова:} Ако је нови сигнал паран, он проверава врх бафера. Ако је на врху такође паран број, долази до интерференције -- оба сигнала се поништавају и нестају из бафера.
	\item \textbf{Непарно чишћење:} Ако је нови сигнал непаран, он редом уклања (брише) све парне сигнале са врха бафера док не наиђе на први непаран сигнал или док се бафер не испразни. Након тога, тај непарни сигнал остаје на врху.
	\item \textbf{Стабилност:} Ако је нови сигнал паран, а на врху бафера је непаран број, он се само додаје на врх без икаквих промена.
\end{itemize}

Твој циљ: Имплементирати методу која прима низ сигнала и враћа финални садржај бафера.

\textbf{Додатни захтеви за решавање задатка:}
\begin{itemize}
	\item Временска сложеност: $O(n)$.
	\item Имплементацију писати користећи одговарајућу структуру података (стек) за ефикасно управљање врхом бафера.
\end{itemize}

\textbf{Пример:}
\begin{itemize}
	\item \textbf{Улаз 1} (несортирани низ целих бројева): \texttt{[3, 4, 6, 5]}
	\item \textbf{Излаз 1:} \texttt{[3, 5]}
	\item Објашњење: \texttt{6} поништава \texttt{4} јер су оба парна; \texttt{5} затим седа на \texttt{3}.
	\item \textbf{Улаз 2} (несортирани низ целих бројева): \texttt{[8, 10, 7, 2, 4, 6]}
	\item \textbf{Излаз 2:} \texttt{[7, 6]}
	\item Објашњење: \texttt{10} поништава \texttt{8}; \texttt{7} улази; \texttt{2} седа на \texttt{7}; \texttt{4} поништава \texttt{2}; \texttt{6} седа на \texttt{7} -> \texttt{[7, 6]}.
\end{itemize}

\sablon{Шаблон другог поправног теста фебруар 2026. група Б}{2025/Februar/GrupaB/AISP2025_PopravniTest2_Test1_GrupaB.linq}
